\section{Introduction}
\label{sec:intro}

%
A key problem in lattice field theory and statistical mechanics is the evaluation of integrals over field configurations, referred to as path integrals. Typically, such integrals are evaluated via a Markov chain Monte Carlo (MCMC) approach: field configurations are sampled from the desired probability distribution, dictated by the action of the theory, using a Markov chain.
%
A significant practical concern is the existence of correlations between configurations in the chain. Critical slowing down~\cite{Wolff:1989wq} refers to the divergence of the associated autocorrelation time as a critical point in parameter space is approached. This behavior drastically increases the computational cost of simulations in these parameter regions~\cite{DelDebbio:2004xh,Meyer:2006ty}. For some models, algorithms have been found which significantly reduce or eliminate this slowing down~\cite{Kandel:1988zz,Bonati:2017woi,PhysRevD.98.054502,SwendsenWang87,Wolff1989,ProkofevSvistunov01,KawashimaHarada04,Bietenholz1995MeronCluster}, enabling efficient simulation. For field theories, a number of methods have been proposed to circumvent critical slowing down by variations of Hybrid Monte Carlo (HMC) techniques~\cite{Ramos:2012bb,Gambhir:2015rha,Cossu:2017eys,Jin2019QNHMC}, multi-scale updating procedures~\cite{Endres:2015yca,Detmold:2016rnh,Detmold:2018zgk}, open boundary conditions or non-orientable manifolds~\cite{Luscher:2012av,Mages:2015scv,Burnier:2017osu}, metadynamics~\cite{Laio:2015era}, and machine learning tools~\cite{Tanaka:2017niz,Shanahan:2018vcv}. In important classes of theories, however, critical slowing down remains limiting; for example, in lattice formulations of Quantum Chromodynamics (QCD, the piece of the Standard Model describing the strong nuclear force) it is a major barrier to simulations at the fine lattice spacings required for precise control of the continuum limit.

Here, a new \emph{flow-based MCMC} approach is proposed and is applied to lattice field generation. The resulting Markov chain has autocorrelation properties that are systematically improvable by an optimization (training) step before sampling.
In this method, samples $z$ are drawn from a simple distribution and then transformed by a change-of-variables (or ``flow'') $\phi = f^{-1}(z)$, resulting in samples $\phi$ with a new effective distribution $\netp$. The mapping $f^{-1}$ is chosen to be efficient to compute, making it easy to draw samples $\phi$, and is optimized within a variational family to produce a distribution $\netp$ close to the desired one.
To guarantee asymptotic exactness of sampling, a Markov chain is constructed using Metropolis-Hastings steps with $\tilde{p}_{f}$ taken as a proposal distribution. Since proposed samples are independent of the previous samples in the chain, the autocorrelation time and acceptance rate are coupled; the autocorrelation time drops to zero as the acceptance rate approaches $1$. This is true even in regions of parameter space where standard algorithms exhibit critical slowing down. Under mild conditions (detailed in Section~\ref{sec:flowMCMC}), this approach is guaranteed to generate samples from the desired probability distribution in the limit of a large number of updates.

This method has several features that make it attractive for the evaluation of path integrals in lattice field theories:
%
\begin{enumerate}
    \item The autocorrelation time of the Markov chain can be systematically decreased by training the model;
    \item
    Each step of the Markov chain requires only the model evaluation and an action computation;
    \item Each update proposal is independent of the previous sample, thus proposals can be generated in parallel and efficiently composed into a Markov chain;
    \item The model is trained using samples produced by the model itself, without the need for existing samples from the desired probability distribution.
\end{enumerate}

Several other machine learning approaches have been applied to MCMC, for statistical mechanics systems, synthetic distributions, and simple lattice quantum field theories. Self-learning Monte Carlo (SLMC) methods have been applied fairly successfully to one- to three-dimensional Ising and fermionic systems. These methods construct, by a variety of techniques, an effective Hamiltonian for a theory that can be more easily sampled than the original Hamiltonian~\cite{Huang2017,Liu2017SLMC,Liu2016SLMCFermion, Nagai2017SLMCCT,Shen2018SLMCDNN}.
The effective Hamiltonian is learned using supervised learning techniques based on training data drawn from a combination of existing MCMC simulations, randomly-mutated samples, and the accelerated Markov chain itself (hence the term ``self-learning'').
Flow-based methods have been used for Monte Carlo sampling in the two-dimensional Ising model~\cite{Li2018}, many-body systems~\cite{NoeBoltzmannGenerators2018}, and
synthetic distributions~\cite{Song2017,Levy2017GenHMC}, and generative adversarial methods have been applied to two-dimensional scalar field theory~\cite{Urban2018,Zhou2018}.

In contrast to these approaches, the method proposed here focuses on directly generating samples from a close approximation to the true distribution in such a way that the exact likelihood of each produced sample is known. The direct generation allows self-learning as in SLMC, and the known likelihood allows use of a Metropolis-Hastings acceptance step to ensure exactness.


The proposed flow-based MCMC algorithm is detailed in Section~\ref{sec:flowMCMC}. A numerical study of its effectiveness in the context of two-dimensional $\phi^4$ theory is presented in Section~\ref{sec:phi4}. Finally, Section~\ref{sec:summary} outlines the further development and scaling of the approach that will be required for applications to theories defined in a larger number of spacetime dimensions and to more complicated field theories such as QCD.
